\scanpage{140}
de $\unsure{\mathrm{Gl}|\mathrm{Sl}}(2,\mathbb Z)$ sur les structures de Torelli, par $g\cdot u=u\circ g^{-1}$ ($u\colon\mathbb Z^2\xrightarrow{\sim}\Pi_E$).

\uline{Corollaire 1} On a un iso. canonique
\[
 \mathscr T_{1,1}\simeq\mathrm{Sl}(2,\mathbb Z) \tag*{(26)}
\]
(où $\mathscr T_{11}=\pi_1(M_{11})$ est calculé relativement au revêtement universel choisi $\mathscr D$ de $M_{1,1}$).
\[
 \text{On a un iso. can}\quad M_{11}\simeq(\deleted{\unsure{\Sigma}}\mathscr D_{\deleted{\unsure{1}}},\mathrm{Sl}(2,\mathbb Z) \tag*{(27)}
\]

Utilisant les grands théorèmes \unsure{marteaux-pilons}, on en déduit

\uline{Corollaire 2} Soit $(\deleted{\unsure{\Sigma|E}}E,a)$ une surface \inserted{ponctuée} orientée compacte connexe de genre $1$, et soit $\mathscr T$ le groupe des classes d'isotopie des automorphismes de $(E,a)$. \deleted{Alors Considérons} \inserted{Alors} l'homomorphisme naturel
\[
 \mathscr T\longrightarrow\mathrm{Aut}_{\mathbb Z}^{+}(H)\qquad\text{où}\quad H=\pi_1(E,a) \tag*{(27)}
\]
est un isomorphisme. \unsure{\S|\deleted{§}}

\uline{Corollaire 3} Soit $\Pi$ un groupe à lacets de type \deleted{\unsure{$g$}} $(1,1)$, $H=\Pi_{\mathrm{ab}}$ (qui est un $\mathbb Z$-module libre de rang $1$), $\mathscr T=\mathrm{Autext}_{\unsure{\mathrm{lac}|\mathrm{loc}}}^+(\Pi)$, \deleted{considérons}\inserted{alors,} l'hom
\[
 \mathscr T\longrightarrow\mathrm{Aut}_{\mathbb Z}^+(H) \tag*{(28)}
\]
est iso.
% slanted
\margin{Il faudrait expliciter l'action \inserted{extérieure} de $\mathrm{Aut}_{\mathbb Z}^+(H)$ sur $\Pi$, \unsure{resp.|puis} l'action $\unsure{-|\text{de}}\ \mathrm{Aut}_{\mathbb Z}^+(H)$ \ldots}
% archive sheet number: 139
